\section{Penyiapan Data dan Pelabelan}

Tahap penyiapan data berfokus pada pembersihan data awal dan pemberian label (\textit{labeling}) untuk pengelompokan data. Label sangat penting untuk melatih algoritma pembelajaran (\textit{supervised learning}) seperti yang dibahas pada Subbab~\ref{sec:machine_learning}. Dalam sistem ini, data penyewa yang telah dikumpulkan akan diseragamkan terlebih dahulu menjadi huruf kecil (lowercasing) dan dihilangkan spasi berlebihnya (stripping). Setelah itu, data akan diberikan nilai berupa angka 0 atau 1 (skor biner). Nilai ini ditentukan secara otomatis dari ada atau tidaknya catatan pelanggaran pada kolom komentar, dengan aturan matematis sebagai berikut.

\[
Label(x) =
\begin{cases}
0, & \text{jika } komentar \in \{ \text{kosong, ``nan'', ``na'', ``null'', ``-'', ``0'', ``tidak ada'', ``none'', spasi} \} \\
1, & \text{lainnya (terdapat catatan pelanggaran)}
\end{cases}
\]

Pendekatan ini memisahkan penyewa ke dalam dua kelompok utama:
\begin{itemize}
    \item \textbf{Kelas Normal (0):} penyewa yang mematuhi aturan dan riwayat catatannya bersih dari pelanggaran.
    \item \textbf{Kelas Tidak Normal (1):} penyewa yang terindikasi melanggar aturan dan menimbulkan kerugian bagi pihak pengelola.
\end{itemize}

Melalui proses penyiapan data dan pelabelan ini, model algoritma dapat menggunakan angka 0 dan 1 tersebut sebagai patokan dasar. Tujuannya adalah agar sistem mampu mengenali pola kebiasaan yang tidak wajar (\textit{anomali}) berdasarkan informasi transaksi penyewa lainnya.